% SHARD-ID: AQM-15-PROJECTIVE-SHEAF-FIELD-DEFINITIONS
% SHARD-TITLE: Projective Sheaf Field Definitions
% SHARD-SUMMARY: Gives self-contained definitions of presheaves, sheaves, global sections, varieties, projective space, and twists.
% SHARD-SUMMARY: Defines the projective replacement for naive Hom(X,V) as a sheaf-section field space with a finite-point pairing.
% SHARD-SUMMARY: Proves the basic vector-space and symplectic-form lemmas needed before the projective-line examples.
% SHARD-KEYWORDS: presheaf, sheaf, projective variety, projective space, twist, arithmetic quantum field

\section{Projective Sheaf Field Definitions}
\label{sec:projective-sheaf-field-definitions}

\subsection{Status and Ground Truth}

This shard is definition-building, not a classification theorem.  It is
source-local where it recalls standard algebraic geometry and checked where it
proves the finite linear-algebra statements.  The local ground truth is the
Stacks Project TeX snapshot registered in
\path{references/algebraic_geometry/SOURCES.md}.  We use:
\begin{itemize}
  \item presheaves and sections:
  \path{stacks_project_sheaves.tex}, lines 73--123;
  \item sheaves: \path{stacks_project_sheaves.tex}, lines 500--525;
  \item Proj and its sheaves:
  \path{stacks_project_constructions.tex}, lines 939--1225;
  \item twists \(\mathcal O(n)\):
  \path{stacks_project_constructions.tex}, lines 1695--1705;
  \item projective space:
  \path{stacks_project_constructions.tex}, lines 2775--2886;
  \item varieties and projective varieties:
  \path{stacks_project_varieties.tex}, lines 53--58 and 4843--4865;
  \item \(H^0(\mathbf P^n,\mathcal O(d))\):
  \path{stacks_project_coherent.tex}, lines 1557--1622.
\end{itemize}

\subsection{Lamport Dependency Skeleton}

\[
\begin{array}{rcl}
D1&:&X\text{ is a topological space.}\\
D2&:&\mathcal F\text{ is a presheaf on }X.\\
D3&:&\mathcal F\text{ is a sheaf if compatible local sections glue uniquely.}\\
D4&:&\Gamma(U,\mathcal F)=\mathcal F(U)=H^0(U,\mathcal F).\\
D5&:&X\text{ is a variety, projective variety, or projective space over }F_q.\\
D6&:&\mathcal O_X(d)\text{ is the }d\text{th twist on a projective space.}\\
D7&:&(V,\omega_V)\text{ is a finite symplectic }F_q\text{-vector space.}\\
D8&:&E(X,\mathcal L,V)=H^0(X,\mathcal L)\otimes_{F_q}V.\\
D9&:&\Omega\text{ is the finite rational-point pairing after choosing evaluations.}\\
L1&:&E(X,\mathcal L,V)\text{ is an }F_q\text{-vector space.}\\
L2&:&\Omega\text{ is bilinear and alternating.}\\
L3&:&E/\operatorname{rad}(E)\text{ is the honest Weyl label space.}
\end{array}
\]

\subsection{Presheaves}

Let \(X\) be a topological space.  A presheaf of sets \(\mathcal F\) on \(X\)
consists of two pieces of data.

First, every open subset \(U\subset X\) is assigned a set \(\mathcal F(U)\).
Its elements are called sections over \(U\).  Second, every inclusion of open
sets \(V\subset U\) is assigned a restriction map
\[
  \rho^U_V:\mathcal F(U)\longrightarrow \mathcal F(V).
\]
The maps must satisfy two identities:
\[
  \rho^U_U=\operatorname{id}_{\mathcal F(U)},\qquad
  \rho^U_W=\rho^V_W\circ\rho^U_V
  \quad\text{when } W\subset V\subset U.
\]
For \(s\in\mathcal F(U)\), write
\[
  s|_V=\rho^U_V(s).
\]
Thus \(s|_W=(s|_V)|_W\).  The notation
\[
  \Gamma(U,\mathcal F)=\mathcal F(U)=H^0(U,\mathcal F)
\]
means exactly the same set of sections over \(U\); the last two symbols do
not add a new construction in degree zero.

For a fixed field \(F_q\), a presheaf of \(F_q\)-vector spaces is the same
data with each \(\mathcal F(U)\) an \(F_q\)-vector space and each restriction
map \(F_q\)-linear.  Equivalently, it is a presheaf of sets whose section sets
carry compatible vector-space operations.

\subsection{Sheaves}

A sheaf is a presheaf satisfying a locality-and-gluing axiom.  Let
\[
  U=\bigcup_{i\in I}U_i
\]
be an open cover.  Suppose \(s_i\in\mathcal F(U_i)\) is a section on each
member of the cover.  The family is compatible if for every pair \(i,j\),
\[
  s_i|_{U_i\cap U_j}=s_j|_{U_i\cap U_j}.
\]
The presheaf \(\mathcal F\) is a sheaf if every compatible family has one and
only one glued section \(s\in\mathcal F(U)\) such that
\[
  s|_{U_i}=s_i\qquad\text{for all }i.
\]
This says precisely: a sheaf is an object whose sections are determined
locally, and whose locally compatible data assemble globally.

\subsection{Varieties, Projective Space, and Twists}

Let \(k\) be a field.  In the Stacks convention, a variety over \(k\) is a
scheme \(X\) over \(k\) such that \(X\) is integral and the structure morphism
\[
  X\longrightarrow \operatorname{Spec}(k)
\]
is separated and of finite type.  A projective variety is a variety for which
this structure morphism is projective; equivalently in the cited Stacks
source, it embeds as a closed subscheme of some \(\mathbf P^n_k\).

Projective \(n\)-space over a ring \(R\) is
\[
  \mathbf P^n_R=\operatorname{Proj}(R[T_0,\ldots,T_n]),
\]
where every \(T_i\) has degree \(1\).  Its standard opens are
\[
  D_+(T_i),\qquad i=0,\ldots,n,
\]
and these cover \(\mathbf P^n_R\).  For \(n=1\), the two standard opens
\[
  D_+(X),\qquad D_+(Y)
\]
are affine lines with coordinates \(Y/X\) and \(X/Y\), respectively.

For \(X=\operatorname{Proj}(S)\), the \(d\)th twist of the structure sheaf is
\[
  \mathcal O_X(d)=\widetilde{S(d)}.
\]
In the special case \(X=\mathbf P^n_R\), the local source computes
\[
  H^0(\mathbf P^n_R,\mathcal O(d))
  =
  \begin{cases}
  R[T_0,\ldots,T_n]_d,& d\geq0,\\
  0,& d<0,
  \end{cases}
\]
where \(R[T_0,\ldots,T_n]_d\) is the degree-\(d\) homogeneous part.  Thus
\[
  H^0(\mathbf P^1_{F_q},\mathcal O(d))
  =
  F_q\text{-span}\{X^{d-i}Y^i:0\leq i\leq d\}
\]
for \(d\geq0\).

\subsection{Why Naive Hom Is Not the Projective Default}

For a finite set \(X\), the previous shard used \(E\leq\operatorname{Map}(X,V)\).
For a projective variety this is too naive if \(\operatorname{Hom}(X,V)\) is
read as regular maps to an affine vector space.  A regular map to affine
space is controlled by global regular functions, and the Stacks source records
that a proper geometrically integral variety has
\[
  H^0(X,\mathcal O_X)=k.
\]
So for connected projective spaces this would see only constants.

The projective replacement is therefore sheaf-theoretic:
\[
  E(X,\mathcal L,V)=H^0(X,\mathcal L)\otimes_{F_q}V,
\]
where \(\mathcal L\) is an explicitly chosen sheaf on \(X\).  The sheaf is part
of the quantum-field datum.  Different choices \(\mathcal L=\mathcal O(d)\)
give different finite field-label spaces.

\subsection{Finite Symplectic Target and Pointwise Form}

Let \(V\) be a finite-dimensional \(F_q\)-vector space.  A symplectic form on
\(V\) is an \(F_q\)-bilinear pairing
\[
  \omega_V:V\times V\to F_q
\]
such that \(\omega_V(v,v)=0\) for all \(v\) and the only vector pairing to
zero with every vector is \(0\).  The standard two-dimensional example is
\[
  V=F_q^2,\qquad
  \omega_V((a,b),(c,d))=bc-ad.
\]

To turn \(E(X,\mathcal L,V)\) into a finite Weyl label space we also need a
finite scalar pairing on \(H^0(X,\mathcal L)\).  Let \(S\subset X(F_q)\) be a
finite set of rational points, let \(w_P\in F_q\), and choose an evaluation
coordinate
\[
  \operatorname{ev}_P:H^0(X,\mathcal L)\to F_q
\]
for every \(P\in S\).  Then define
\[
  B(s,t)=\sum_{P\in S}w_P\,\operatorname{ev}_P(s)\operatorname{ev}_P(t).
\]
This evaluation choice is automatic for honest scalar functions, but for
sections of a nontrivial line bundle it is a convention: one has chosen local
trivializations or fixed projective representatives at the points.

For pure tensors set
\[
  \Omega(s\otimes v,t\otimes w)=B(s,t)\omega_V(v,w),
\]
and extend bilinearly to all of \(E(X,\mathcal L,V)\).

\paragraph{Lemma \(L1\).}
\(E(X,\mathcal L,V)\) is an \(F_q\)-vector space.

\emph{Proof.}
The source definition of a sheaf of vector spaces gives
\(H^0(X,\mathcal L)=\mathcal L(X)\) as an \(F_q\)-vector space.  The tensor
product of two \(F_q\)-vector spaces is an \(F_q\)-vector space, with addition
and scalar multiplication inherited from the tensor construction.  Hence
\(H^0(X,\mathcal L)\otimes_{F_q}V\) is an \(F_q\)-vector space.

\paragraph{Lemma \(L2\).}
\(\Omega\) is \(F_q\)-bilinear and alternating.

\emph{Proof.}
The map \(B\) is bilinear because it is a sum of products of linear evaluation
maps.  The form \(\omega_V\) is bilinear by definition.  Therefore their
tensor product pairing is bilinear on pure tensors and hence bilinear after
extension.  Also
\[
  \Omega(s\otimes v,s\otimes v)=B(s,s)\omega_V(v,v)=0,
\]
and the same alternating identity follows for arbitrary sums by bilinearity
and the alternating property of \(\omega_V\).

\paragraph{Lemma \(L3\).}
Let
\[
  \operatorname{rad}(E)=
  \{\phi\in E:\Omega(\phi,\psi)=0\text{ for every }\psi\in E\}.
\]
Then \(E_{\rm red}=E/\operatorname{rad}(E)\) has a nondegenerate alternating
form induced by \(\Omega\).

\emph{Proof.}
If \(r\in\operatorname{rad}(E)\), then
\(\Omega(\phi+r,\psi)=\Omega(\phi,\psi)\), so the form descends to the
quotient.  If a quotient class \([\phi]\) pairs to zero with every quotient
class, then \(\phi\) pairs to zero with every element of \(E\), hence
\(\phi\in\operatorname{rad}(E)\) and \([\phi]=0\).  Thus the descended form is
nondegenerate.

The Weyl-Heisenberg field displacement algebra is built from
\(E_{\rm red}\), not from unreduced \(E\), exactly as in
Shard \(AQM\)-14.
